\documentclass[onecolumn, draftclsnofoot,10pt, compsoc]{IEEEtran}
\usepackage{graphicx}
\usepackage{url}
\usepackage{setspace}

\usepackage{geometry}
\geometry{textheight=9.5in, textwidth=7in}

\usepackage{graphicx}
\graphicspath{{images/}}

% 1. Fill in these details
\def \CapstoneTeamName{		Team Wombat}
\def \CapstoneTeamNumber{		15}
\def \GroupMemberOne{			Victor Li}
\def \GroupMemberTwo{			Ryan Crane}
\def \GroupMemberThree{			Nicholas Wong}
\def \CapstoneProjectName{		Axolotl}
\def \CapstoneSponsorCompany{	}
\def \CapstoneSponsorPerson{		Kevin McGrath}

% 2. Uncomment the appropriate line below so that the document type works
\def \DocType{		%Problem Statement
				Requirements Document
				%Technology Review
				%Design Document
				%Progress Report
				}
			
\newcommand{\NameSigPair}[1]{\par
\makebox[2.75in][r]{#1} \hfil 	\makebox[3.25in]{\makebox[2.25in]{\hrulefill} \hfill		\makebox[.75in]{\hrulefill}}
\par\vspace{-12pt} \textit{\tiny\noindent
\makebox[2.75in]{} \hfil		\makebox[3.25in]{\makebox[2.25in][r]{Signature} \hfill	\makebox[.75in][r]{Date}}}}
% 3. If the document is not to be signed, uncomment the RENEWcommand below
%\renewcommand{\NameSigPair}[1]{#1}

%%%%%%%%%%%%%%%%%%%%%%%%%%%%%%%%%%%%%%%
\begin{document}
\begin{titlepage}
    \pagenumbering{gobble}
    \begin{singlespace}
    	\includegraphics[height=4cm]{coe_v_spot1}
        \hfill 
        % 4. If you have a logo, use this includegraphics command to put it on the coversheet.
        %\includegraphics[height=4cm]{CompanyLogo}   
        \par\vspace{.2in}
        \centering
        \scshape{
            \huge CS Capstone \DocType \par
            {\large\today}\par
            \vspace{.5in}
            \textbf{\Huge\CapstoneProjectName}\par
            \vfill
            {\large Prepared for}\par
            \Huge \CapstoneSponsorCompany\par
            \vspace{5pt}
            {\Large\NameSigPair{\CapstoneSponsorPerson}\par}
            {\large Prepared by }\par
            Group\CapstoneTeamNumber\par
            % 5. comment out the line below this one if you do not wish to name your team
            \CapstoneTeamName\par 
            \vspace{5pt}
            {\Large
                \NameSigPair{\GroupMemberOne}\par
                \NameSigPair{\GroupMemberTwo}\par
                \NameSigPair{\GroupMemberThree}\par
            }
            \vspace{20pt}
        }
        %\begin{abstract}
        % 6. Fill in your abstract    
   % Our project requires the development of software to run a car infotainment system alongside a black box data log. Our software will run on an Nvidia Jetson TX2 and should be able to execute basic media and navigation tasks while also logging data from a variety of sensors. The goal is to allow people to modernize older cars with an array of fresh sensors and an interactive software interface. To accomplish this goal, we will pull vehicle data from a car?s OBDII (On Board Diagnostics II) port as well as our own sensors wired into the TX2, and feed that data to software that aggregate that data into a log. 
      %  \end{abstract}     
    \end{singlespace}
\end{titlepage}
\newpage
\pagenumbering{arabic}
\tableofcontents
% 7. uncomment this (if applicable). Consider adding a page break.
%\listoffigures
%\listoftables
\clearpage

% 8. now you write!
\section{Introduction}
\subsection{Purpose}
The purpose of this software requirements specification (SRS) document is to outline and detail the capabilities of the NVIDIA Jetson TX2 infotainment and black box our group will develop, known henceforth as the Axolotl Infotainment System and Axolotl OS. Doing so will enable us to describe the requirements of the Axolotl Infotainment System and Axolotl OS such that we and our client will have a detailed understanding of the form factor and capabilities of the deliverable system we will develop. The intended audience for this SRS includes our client, the CS Capstone Instructors, and our group.\par

\subsection{Scope}
Our project entails the development of an infotainment system and black box that can be divided into two products: the Axolotl and the Axolotl Software. The Axolotl will connect vehicle sensors, controllers, receivers, and a touchscreen to a NVIDIA Jetson TX2 computer in a package that can be installed in a vehicle. The Axolotl Software runs on the Axolotl and provides users with media playback, navigation, and vehicle data logging capabilities.\par


\section{Definitions}
\begin{itemize}
	\item NVIDIA Jetson TX2: A versatile, efficient, and high-performance computer made by NVIDIA to be used in robots, drones, and smart cameras.

	\item OBD-II: OBD-II refers to On-Board Diagnostics II, a standardized connector installed in all automobiles sold in the United States since January 1st, 1996. Devices connected via a car’s OBD-II port can read the vehicle’s sensor data on the fly.

	\item LIDAR: LIDAR refers to Light Detection and Ranging, a method of using laser pulses to determine the 3D properties of a faraway object.

	\item AHRS: AHRS refers to Attitude, Heading, and Reference System, a system used in modern aircraft to determine and display roll, pitch, and yaw.

	\item Infotainment: A portmanteau of “information” and “entertainment”. When we use the term “infotainment”, we are referring to the center console touchscreen that gives drivers access to information and media in modern cars.

	\item RDS: RDS refers to Radio Data System, a method of transmitting the current track information of an FM radio broadcast.

	\item UPS: UPS refers to Uninterruptible Power Supply, an auxiliary power source that enables a device to function for a limited time if its main power source is unavailable.

\end{itemize}

\section{Overall Description}
\subsection{Product Perspective}
The Axolotl Infotainment System is comprised of the Axolotl Head Unit and the Axolotl Software.
The Axolotl Head Unit consists of all of the necessary sensors, receivers, and controller hardware connected physically and wirelessly with the NVIDIA Jetson TX2 system. The Axolotl is designed to be integrated into a car to either provide or replace an in-car infotainment system. Users will not interact directly with the sensors, receivers, or controllers.\par

The Axolotl Software will be installed on the Axolotl Head Unit’s TX2 unit and directly interface with the user. The software will operate on Linux, as it is the base operating system installed on the TX2. The Axolotl Software will interact with the hardware and provide users with the ability to control media playback, conduct mapping and routing with navigation, and also exert limited control over system settings.\par

\subsection{Product Functions}
The main functions of the Axolotl are:
\begin{itemize}
	\item	The Axolotl will allow users to play media from multiple sources including: USB, Bluetooth, Auxiliary, and FM.

	\item	It will also offer navigation with destination entry, mapping, and offline capabilities. 
	\item	The black box portion of the Axolotl logs output from a dashcam and all sensors tied into a car’s OBD-II port. Users are able to download the black box data to a storage device or clear all black box data. 
	\item	The Axolotl display will also switch to the backup camera whenever user is reversing the car.
	
\end{itemize}
\section{Appendix... maybe}
\includegraphics{universe}

\end{document}
